% ============================================================================
\clearpage
\section*{v4.10.0-r01 Source-Only: The Full Einstein Equation of W1 and the Sign of the Gravitational Action}
\addcontentsline{toc}{section}{v4.10.0-r01 Source-Only: The Full Einstein Equation of W1 and the Sign of the Gravitational Action}
\label{sec:v4100-r01-einstein-sign}
% ============================================================================

\paragraph{Purpose.}
Layer v4.9.0-r86 derived the field equations of the joint action W1 (\S\ref{sec:v490-r81-joint-action-w1}) from their definitions up to
the full non-linear Einstein equation and the Bianchi identity (\S\ref{sec:v490-r86-wolfram-checks}).  This first layer after the
v4.10.0 release closes that gap with one pre-registered test (H90, not blind) and records what it found on the way: W1 as written has the
wrong sign of the gravitational action.  The author decided on 6 November 2026 to correct it.  The layer states the corrected action and
lists which earlier statements change sign.  Earlier layers are not rewritten.  The notes are
\texttt{notes/\allowbreak research/\allowbreak v4\_10\_0\_r01/}.

\noindent\textbf{Source-only status.}
This v4.10.0-r01 layer is source-only.  It does not build a PDF, does not assign a new DOI,
and does not authorize a public release.  The v4.10.0 Version DOI
\href{https://doi.org/10.5281/zenodo.23991357}{10.5281/zenodo.23991357} and the Concept DOI
\href{https://doi.org/10.5281/zenodo.19210728}{10.5281/zenodo.19210728} are unchanged by it.

\subsection*{v4.10.0-r01 results}
\label{subsec:v4100-r01-results}

\paragraph{Theoremlet v4.10.0-r01.1 (H90: the Einstein--Cartan equations of W1, non-linear).}
Pre-registered (commit \texttt{e5d03d8}, not blind).  Conventions as in r86: \(\eta=\operatorname{diag}(1,-1,-1,-1)\),
\(D_\mu\psi=\partial_\mu\psi+\tfrac12\omega_{\mu ab}\Sigma^{ab}\psi+iqJA_\mu\psi\),
\(R^{ab}{}_{\mu\nu}=\partial_\mu\omega_\nu{}^{ab}-\partial_\nu\omega_\mu{}^{ab}+\omega_\mu{}^a{}_c\,\omega_\nu{}^{cb}-\omega_\nu{}^a{}_c\,\omega_\mu{}^{cb}\), and the gravitational action
written as \(S_{\rm grav}=\tfrac{\sigma_g}{2\kappa}\int e\,R\) (W1 as written in r81.1: \(\sigma_g=+1\)).  At a generic point of three random curved
backgrounds (vierbein a random polynomial of degree three, not diagonal) all fields were carried as jets to third order: three Weyl
fields with \(J=(+1,-1,+1)\), a real potential \(A\), the vierbein.  The spinor jets solve the Weyl equation with gauge coupling and
four-fermion term at the point and to first derivative, \(A\) the Maxwell equation at the point.  Both signs \(\sigma_g=\pm1\) were computed.
\begin{enumerate}[leftmargin=2em]
\item[(a)] Cartan equation.  The spin connection \(\tilde\omega\) of the tetrad postulate solves the full \(\omega\)-equation (with its derivative
term) without spinors; with spinors \(\omega=\tilde\omega+K\) solves it, with \(K_{cab}=-\tfrac{\sigma_g}{4}\kappa\,\varepsilon_{cabd}\,j^d\) totally
antisymmetric (residuals \(\le2\times10^{-14}\)).
\item[(b)] Einstein--Cartan equation.  The variation with respect to the vierbein at fixed \(\omega=\tilde\omega+K\) is symmetric on the field
equations and equals a multiple of \(\tilde G_{\mu\nu}+\sigma_g\kappa\,T^{\rm eff}_{\mu\nu}\) (residual \(10^{-15}\)), with
\[
T^{\rm eff}_{\mu\nu}=T^{\rm Bel}_{\mu\nu}-g_{\mu\nu}\mathcal L_W+T^{\rm Maxwell}_{\mu\nu}-g_{\mu\nu}\,\lambda\,j\cdot j,
\qquad \lambda=-\sigma_g\,\tfrac{3\kappa}{16}.
\]
Hence \(\tilde G=-\sigma_g\,\kappa\,T^{\rm eff}\), fully non-linear.  The \(K\) term of the Einstein--Cartan Weyl equation equals
\(2\lambda\,j_a\bar\sigma^a\psi\).
\item[(c)] Compatibility.  \(\tilde\nabla^\mu T^{\rm eff}_{\mu\nu}=0\) on the field equations (relative \(4\times10^{-15}\)); \(\tilde\nabla^\mu\tilde G_{\mu\nu}=0\) as a
control.  Negative controls: without the four-fermion term the residuals rise to \(0.12\)--\(0.19\) (Einstein) and \(0.25\)--\(0.81\)
(divergence); with spinors off the field equations to \(0.16\)--\(0.19\) and \(0.92\)--\(1.05\).
\item[(d)] Sign.  For W1 as written (\(\sigma_g=+1\)) the factor is \(-1\) to \(4\times10^{-15}\): \(\tilde G=-\kappa T^{\rm eff}\).  A second system
(Wolfram) gives \(G_{00}=2\nabla^2\Phi\) for \(g=\operatorname{diag}(1+2\Phi,-(1-2\Phi)\mathbb 1)\) in the same curvature convention, so
\(\tilde G=+\kappa T\) is attraction and \(\tilde G=-\kappa T\) repulsion.  W1 as written repels positive energy.
\end{enumerate}
The Hehl--Datta control of r81.1(b) compared only the modulus \(3/16\); the sign had not been checked.

\paragraph{Decision v4.10.0-r01.2 (author, 6 November 2026: the corrected action).}
W1 is corrected to
\[
S=\int e\,\Bigl[-\tfrac1{2\kappa}R(e,\omega)+\tfrac12\bigl(\Psi^\dagger i\bar\sigma^\mu D_\mu\Psi+\mathrm{c.c.}\bigr)-\tfrac14F^2\Bigr]
\]
in the signature \((+,-,-,-)\) (as in Landau--Lifshitz).  The field equations become \(\tilde G=\kappa\,T^{\rm eff}\) with
\(\mathcal L_{4f}=+\tfrac3{16}\kappa\,j\cdot j\), \(T^{(4f)}_{\mu\nu}=-\tfrac{3\kappa}{16}\,j^2g_{\mu\nu}\), Weyl term \(+\tfrac{3\kappa}{8}j_a\bar\sigma^a\psi\),
\(K_{cab}=+\tfrac\kappa4\varepsilon_{cabd}j^d\).

\paragraph{Correction v4.10.0-r01.3 (statements that change sign).}
Moduli and forms stay; the signs change as follows.
\begin{enumerate}[leftmargin=2em]
\item[(a)] r81.1(b) and r86.2(b), (c): \(\mathcal L_{\rm eff}=+\tfrac3{16}\kappa\,j\cdot j\) and \(+\tfrac{3\kappa}{8}j_a\bar\sigma^a\psi\) in the Weyl equation (were
\(-\)); correct for the action as written, not for the physical one.
\item[(b)] r86.3(a), (b): \(T^{(4f)}=-\tfrac{3\kappa}{16}j^2g\) (\(\Lambda\) form, negative); the homogeneous state of two opposite spins has
\(\rho=p=-\tfrac{3\kappa}{16}n^2<0\), the negative spin energy density known from Einstein--Cartan theory.
\item[(c)] Correction r86.4 is withdrawn in its sign part: the original reading of the W1 notes, \(\rho\propto-\kappa n^2\), holds for the physical
sign.  ``Not like \(\Lambda\)'' for the whole tensor and r81.1(c) (no vacuum term) stand.
\item[(d)] r82.4: the pairing channel of the torsion interaction is repulsive, \(H_{\rm int}\supset+\tfrac34\kappa\,\Phi^\dagger\Phi\).  No condensate arises
through it; the \(\nu^c\) Majorana mass near \(M_P\) and the seesaw estimate of r82.4 fall.
\item[(e)] Not affected (review in \texttt{llm/\allowbreak Durchsicht\_\allowbreak Vorzeichen\_\allowbreak kappa\_\allowbreak r81\_\allowbreak r86.md}): H69--H73 (source side only),
H74 (no gravity), H75 (computed with Newtonian attraction), H76 (rate \(\propto\kappa^4\)), H77--H81 and H84--H86 (no gravitational coupling), H82
and H83 (relative weights and phases of the Einstein--Cartan term), H48 (Hehl--Datta sign taken from the literature).
\end{enumerate}

\paragraph{Proof and measurement sources.}
Script \texttt{k205\_\allowbreak w1\_\allowbreak einstein\_\allowbreak nichtlinear.py} (Python, jets, output \texttt{k205\_\allowbreak out.json}) and Wolfram code
\texttt{k205b\_\allowbreak schwachfeld.wl}; notes in \texttt{notes/\allowbreak research/\allowbreak v4\_10\_0\_r01/}; pre-registration and result H90 in
\texttt{notes/\allowbreak research/\allowbreak predictions/}.  \(\square\)

\paragraph{Consequence.}
The field equations of W1 are now derived independently in full: Cartan equation, elimination of the torsion, Einstein--Cartan equation, its
Riemannian form with the four-fermion term, and its compatibility with the Bianchi identity.  With the corrected sign the four-fermion
part has the form of a negative vacuum term with a height \(\propto j^2\); the whole tensor of a homogeneous state is stiff matter with
negative energy, not \(\Lambda\).  W1 still produces no vacuum term.  OP5 is not closed.

\subsection*{v4.10.0-r01 status in the three frames}
\label{subsec:v4100-r01-frames}

\begin{itemize}[leftmargin=2em]
\item r01.1--r01.3 concern the matter part of W1 in frame M.  The tachyon image under the second reading (r84) multiplies the action by the
weights \(a\), \(b\); the sign correction is common to both and does not change their ratio.  The photon frame P is the degenerate fixed point;
no frame-free statement is made beyond that.
\item r01.3(d) concerns the pairing of two left-handed fields in frame M; under the second reading the image term is \(-\mathcal L_M\) (r84), so the
interaction keeps its sign relative to the kinetic term.
\end{itemize}

\subsection*{v4.10.0-r01 register entries}
\label{subsec:v4100-r01-register}

\begin{center}
\small
\begin{tabular}{@{}>{\raggedright\arraybackslash}p{0.34\textwidth}>{\raggedright\arraybackslash}p{0.58\textwidth}@{}}
\toprule
statement & status after v4.10.0-r01\\
\midrule
full Einstein equation of W1, Bianchi & derived and compatible (H90)\\
sign of the gravitational action & corrected to \(-\tfrac1{2\kappa}\int eR\) (author decision)\\
four-fermion term & \(+\tfrac3{16}\kappa\,j\cdot j\); \(T^{(4f)}\) negative \(\Lambda\) form\\
homogeneous spin state & \(\rho=p=-\tfrac{3\kappa}{16}n^2<0\)\\
torsion pairing channel (r82.4) & repulsive; Planck-scale \(\nu^c\) condensate withdrawn\\
W1 vacuum term & still absent\\
\bottomrule
\end{tabular}
\end{center}

\subsection*{v4.10.0-r01 ledger and guard update}
\label{subsec:v4100-r01-ledger-guard}

\[
R_{\rm chk}^{v4100\text{-}r01}=(1,1,1)
\]
for (Einstein--Cartan equations derived; compatibility shown; sign corrected by author decision).

The layer spends no guard:
\[
R_{\rm glob}^{(v4100\text{-}r01,+)}=(0,1,1,1,1),
\qquad
\Gamma_{\rm guard}^{v4100\text{-}r01}=(1,1,1,1,1),
\qquad
\beta_4=1/256\ \text{unspent}.
\]
It does not imply exact finite \(C5\) zero closure, exact all-readout terminality,
\(C6/T2\), a global photon-clock/RH theorem, or global all-readout closure.
